This section fixes one canonical main-text answer to the question ``What is $\Phi_H(s)$?''
It does not claim this choice is already validated.
Its role is to state the exact object and list the theorem burdens required before RH-reduction claims become theorem-level statements.

\subsection*{Canonical main-text choice}
\placeholder{The operator construction and regularization theorem package are not yet complete.}
The manuscript fixes one functional in the main flow:
\[
\Phi_H(s):=\Tr_{\mathrm{reg}}(H-sI)^{-1}-\Tr_{\mathrm{reg}}(H-(1-s)I)^{-1},
\]
defined on
\[
\mathcal{A}:=
\left\{
s\in\C:
\Tr_{\mathrm{reg}}(H-sI)^{-1},\Tr_{\mathrm{reg}}(H-(1-s)I)^{-1}\ \text{exist and are finite}
\right\}.
\]
The associated predicate is
\[
\text{Bal}_H(s)\iff \Phi_H(s)=0 \qquad (s\in\mathcal{A}).
\]
Interpretation: $\Phi_H(s)$ is the regularized resolvent reflection defect of $H$ at the pair $(s,1-s)$.

\subsection*{Minimal required structure for the $H$-layer}
The manuscript does not need every possible operator-theoretic feature.
It does need the following \emph{minimal interface} to keep later arguments well-typed and auditable.

\paragraph{(M1) Ambient space and typing data.}
A specified complex topological vector space (typically a Hilbert or Banach model) $\mathcal{X}$, together with clear codomain conventions for all derived maps.
Every use of $H$ must state its typing context explicitly.

\paragraph{(M2) Domain and well-defined action.}
A dense or otherwise explicitly justified domain $\mathcal{D}(H)\subseteq\mathcal{X}$ (as required by the chosen model), and a map
\[
H:\mathcal{D}(H)\to\mathcal{X}
\]
with proof that this action is well defined under stated hypotheses.
If closure/closability is needed by subsequent lemmas, that requirement must be stated and proved where used.

\paragraph{(M3) Symmetry-compatibility package.}
A precise compatibility statement linking the operator layer to the reflection involution $\sigma(s)=1-s$ from \texttt{definitions.tex}.
This must be explicit and non-heuristic.
Informal symmetry pictures are admissible only as motivation, not as substitute for this package.

\paragraph{(M4) Operator-regularity assumptions (structural side).}
The exact regularity required for operator statements: continuity/measurability/analytic dependence of operator-valued maps, domain persistence, and closure/graph conditions needed for later lemmas.
Each assumption must be tagged as either proved in-manuscript or imported as an explicit hypothesis.

\paragraph{(M5) Dependency traceability.}
A lemma-level dependency map from operator properties to later claims (especially Targets~A and~B).
No RH-relevant step should rely on an unlisted implicit property of $H$.

\subsection*{Required theorem package for this explicit \texorpdfstring{$\Phi_H$}{Phi_H}}
The definition above is explicit; the open work is to prove it is mathematically sound and suitable for Targets~A/B.

\paragraph{(P1) Non-circular construction.}
The operator model and regularization map $\Tr_{\mathrm{reg}}$ must be defined without using the zero set of $\zeta$.

\paragraph{(P2) Domain of evaluation (admissible class).}
The declared set $\mathcal{A}$ must be proved nonempty and explicitly characterized.
If singular points are removed, the exclusion mechanism must be stated.

\paragraph{(P3) Well-posedness and regularity.}
For $s\in\mathcal{A}$, prove existence and finiteness of both regularized traces and of their difference.
State the continuity/analytic properties of $s\mapsto\Phi_H(s)$ used later.

\paragraph{(P4) Reflection behavior.}
The explicit formula gives
\[
\Phi_H(1-s)=-\Phi_H(s)
\]
whenever both sides are defined.
A theorem-level proof must show that the admissible class is stable under $s\mapsto 1-s$ and that the regularization is compatible with this identity.

\paragraph{(P5) Bridge to zero claims.}
A precise theorem target must explain how
\[
\Phi_H(s)=0
\]
for this explicit resolvent-defect functional maps to nontrivial zero information.
The bridge should separate:
\begin{enumerate}
  \item the forward direction to be proved,
  \item the converse direction to be proved,
  \item side conditions preventing spurious balanced points.
\end{enumerate}

\subsection*{Canonical exclusion target supported by this interface}
Let
\[
\mathcal{A}_{\mathrm{strip}}:=\{s\in\mathcal{A}:0<\RePart(s)<1\},
\qquad
\mathcal{B}:=\{s\in\mathcal{A}_{\mathrm{strip}}:\Phi_H(s)=0\}.
\]
The main-text exclusion burden is the single subset statement
\[
\mathcal{B}\subseteq\{s\in\mathcal{A}_{\mathrm{strip}}:\RePart(s)=\tfrac12\}.
\]
This section supports that statement by listing operator and analytic prerequisites.
It does not claim those prerequisites are proved.

\subsection*{Open support packages (dependency statements, not proved results)}
\paragraph{Operator-side package (OP).}
The following remain open theorem targets: well-defined operator action and closure properties, existence of the regularized resolvent traces, regularization invariance/consistency, and dependency-locked usage of these results across Targets~A/B.

\paragraph{Analytic-control package (AC).}
The following remain open theorem targets: strip-wide bound control for the explicit $\Phi_H$, legitimacy of limit/interchange operations for regularized traces, continuation control from local to strip-wide regimes, and exclusion of hidden off-line branches of $\mathcal{B}$ caused by analytic degeneration.

\paragraph{Status note.}
OP and AC are recorded as open bridge classes for audit.
They specify proof burden only.
They do not provide a proof of Target~A or Target~B.

\subsection*{Why this explicit choice is kept in the main flow}
The conditional reduction step in \texttt{main-results.tex} depends on theorem-level statements about the specific functional above.
Without a rigorous foundation for this explicit $\Phi_H$:
\begin{enumerate}
  \item $\Phi_H(s)=0$ is not yet a theorem-level condition on a proved domain,
  \item the zero--balance correspondence target lacks a justified bridge,
  \item the canonical exclusion target lacks a proved analytic object to control.
\end{enumerate}
Therefore this explicit $\Phi_H$ choice is a prerequisite object for the manuscript's proposed RH reduction route.

\begin{heuristic}
Geometric or cancellation intuition can guide candidate constructions of $H$ and $\Phi_H$.
However, such intuition remains heuristic until converted into explicit definitions, hypotheses, and proved lemmas matching the interface above.
\end{heuristic}
